
%COVID cancelled standardized testing in the UK and lead to record number of  top grades getting issues for two years straight, and overstreching universities capacities.
How would teachers and schools grade students differently compared to a standardized grading system where teachers have no discretion?  The COVID 19 pandemic canceled the standardized exams in the UK for 2 consecutive years, and students were permitted to use their teachers' predicted grades as a replacement for the grades. The grade provided by the teachers determined the placement outcomes of the students, as the UK higher education system employs an admission system that determines the success of an application by the applicant grade in the standardized exam (e.g. A-levels).  The policy lead to a 12 percentage point increase in the number of top grades provided in 2020, and a even larger of top grades issued in 2021 albeit the disruption in learning for the 2021 cohort. Consequently, top universities were crowded as they over-accepted students than their capacity, and lower ranked universities suffered from student shortage. This incident underscores how a change in grading system significantly changed the grade distribution with far-reaching consequence on the higher education market.

% Why this is economics? Because its localization of information.
The complete transfer of grading discretion from a centralized examiner to the local teachers triggered a new allocation of grades with far reaching influence on the students subsequent educational and labor market outcomes. The relationship between grades and student characteristics (academic ability or demographics) potentially changed by the teacher's choice to assign grades to students varying characteristics, and also through the grading policies that the school adopted for assigning grades. Previous literature documents the differences across grading patterns between standardized and non-standardized testing. \cite{hanna2012} and \cite{lavy2018} documents how teachers assign better grades to students based on their socioeconomic status as well as their demographics. \cite{figlio2006}  finds institutional differences in degrees of miss-reporting their student's academic competence. \cite{10.1257/app.20200525} find evidence of persistence of grade inflation over time.  However, none of the papers document how grading policies across institutions varies across the entire national secondary school system as well their importance relative to grade inflation due to student's characteristics. Additionally, no papers examine the dynamic persistence in grade inflation practices overtime as well as their granular compositions by smaller subgroups such as subject groups. 

% Why this is economics? Grades determines the labor market outcomes for university graduates. 
The changes in grading practices could potentially grant access for under-represented student groups to secure a placement at selective universities.  \cite{Chetty2023} \cite{Arcidiacono2005}\cite{Epple2005} discusses different admission rules, such as affirmative actions or application assistance, changes the student composition at universities. However, they do not talk about non-standardized grading, but estimate student composition with a combination of standardized tests (e.g. SAT), high-school reports, and admission assistance. 

% How schools inflate differently. Static patterns as well as their dynamics.
I first test whether schools systematically differed in their tendency to assign higher grades based on their attributes. I examine the association of larger leniency towards assigning higher grades with the school type (private school or public schools), and also with their institutional quality. I also test whether the grading patterns differed between students taking different subjects, which I particularly categorize by the extent of quantitative reasoning in their content. Quantitative contents may prevent teachers to over-report their students grades as the past scores of the students from past quiz had already ranked students which facilitated the teachers to predict grades. I estimate the grading patterns while taking into account the student's academic and socio-economic background, as their attending institutions tendency to assign higher grades as well as their sorting behavior.

% I use this case to show how grade distribution changed after the swith, by isolating changes at the instiuttuional and individual level.
I then examine how grades were assigned differently across students with varying demographic and academic characteristics. Namely, I examine whether the extent of the grade inflation was higher for students with higher academic and economic affluence, measured by students performance in another standardized national exam and the students parental income, measured by their parents occupation class. I also compare the relative importance of the students attending schools grading policy and the students' characteristics by the students rank within the student population undertaking the A-levels in a given year. 

% Persistance and how importance change overtime.
I then examine the dynamic patterns of grade inflation within schools during the two consecutive years of the pandemic. I test whether schools changed their grading policies in the second year, while conditioning the effect by look at extent of grade inflation in the initial year. Furthermore, I examine how the distribution of grade inflation across subject groups within a school, connect to the aggregate grade inflation level in the consecutive year, as well as the inflation levels at each subject group. 

% I predict how these grades changed the overall assignment of grades in the student population as well as their consequence in the higher-education market.
Lastly, I provide predictions on how the two grading schemes lead to a different composition of students across universities. In the UK, universities control the number of students they admit by setting a uniform grade requirement that applicants must pass to secure a spot. Using a structural model with endogenous university application choices, with fixed number of seats at universities, I examine how the correlation between university rank and mean demographics of students differs across standardized and non-standardized exams. I document how the share of students from socioeconomically disadvantaged backgrounds at prestigious university institutions differ across the grading regimes. The exercise assesses the extent of the grade changes triggered by the regime change into improvements into placements into better universities, as well as the demographics of students that improves their placement outcomes and those that descend in their placement rank.

\vspace{1em}

I estimate the changes in grades assignments using a binary choice model that distinguishes between student attributes that are common across all schools (e.g., prior academic performance, student demographics) and those specific to schools (e.g. school fixed effects). I measure the changes in grades by the tendency for a given student to receive a top grade. The model estimates the increments in the probability that a student will receive a top grade by comparing students in the non-pandemic years to students in previous cohorts with similar academic standings. To isolate the concentration of grade shifts for each attribute of students, the treatment variable is interacted with a rich set of student demographic measurements, as well as the school indicator. I also include those controls into the regression so that students in the pandemic years are matched with students who took the standardized exam in the non-pandemic years with academically and demographically similar attributes. 

I identify the treatment effects through an interrupted time series (ITS) design that compares the probability of receiving an A or A* across cohorts before and after the grading policy change. The concerns for identification are 1. The treatment effect may include underlying trends in the overall grade distribution, and 2. The treatment effect may not distinguish between changes in the grading method and disruption in student learning due to the closure of schools caused by the pandemic. For the first point, I exploit the institutional grading rule in the UK, in which the overall A-level grades followed a fixed distribution during the years leading to the pandemic. Examiners of the A level set grade boundaries to regulate grade inflation. I tested this property of the A-level distribution through multiple stationarity tests, showing that the average probability of receiving top grades remains stable over years when standardized grading was applied. For the second point, I provide evidence that the learning loss for the 2020 year cohort was small, as school closures were announced a month before the national examination date, and teachers were already informed about their students' academic competence through their teaching experience during the 2 years of the student learning period. [I provide national statistics of mock exams that were taken during the short span between the cancellation of the exams and the teachers' grading date, and show the non-existing evidence of learning disruption for the 2020 year cohort.]

I summarize the within-/between effect using an Oaxaca-Blinder-type decomposition that separates the two effects. I decompose the changes in the mean levels of students' chances of receiving a top grade into within-school effect, which keeps contribution of the school fixed effect as their baseline levels but allowing the student characteristic related factors to change their contribution (e.g. race, gender, previous academic achievements, etc.). I calculate the effect between schools similarly by shutting down the contribution of student characteristics. I also decompose the variance of the individual treatment effect, and decompose the variance into the extent of heterogeneity in both within/between school effects. 

To connect the grade changes to shifts into different university tiers, I develop a structural model of student application and university admission. The model estimates student preferences on universities while controlling for the universities' selectivity levels and the students' aggregate score of admission grades. The models maps grade distributions into the applicant pool for each university through the students' application portfolio. Through a styled assignment rule, in which spots in universities are taken by students from highest achievers, I estimate the full distribution of students across universities in the UK. I provide how the correlation between university selectivity levels and mean levels of the incoming student demographics, such as race and gender, differ between the standardized and non-standardized exam regimes.


\vspace{1em}

I find substantial heterogeneity in how schools varied in their tendency to loosen their grading policies, thereby granting students with higher grades than in the standardized exams. I find that the extent to which they reduced their contribution varied by school quality and their types. School with better quality, measured by their past cohorts exam performances at the A level, were associated with a lower degree of grade inflation, measured by increments in student's chances of receiving top grades, whereas schools with lower tendency to produce a top grade scoring student were associated with higher degrees of grade inflation.  I also find varying grade inflation between school types, with private institutions exhibiting the largest upward shifts. Students in independent schools were 7 percentage points more likely to receive top grades than those attending three-year secondary schools (i.e., sixth-form colleges), which experienced the smallest increases. 

I find that most schools upward revised their grading policies in the second year, thereby loosening their grading standards and contributing to the increase in top grade issued to students. Private schools displayed the largest degree of upward revision in their grading policy terms. The persistence in grading policies were stronger for non-quantitative subjects, with subject contents with stronger discretion of the teachers assessment of the teacher, while quantitative subjects showed a smaller positive persistence in the schools degree of inflating grades. I also show how schools with higher grade inflation in quantitative subjects were show were often related with a larger degree of grade inflation in non-quantitative subject groups in the initial year. 

I find that within a subject course, the assignment of top grades was concentrated among students with academically and socioeconomic advantaged backgrounds. Students in the least economically deprived regions were 3 percentage points more likely to receive an A or A* compared to students in the most deprived regions. In addition, top grades were more likely to be assigned to students with higher scores in the previous academic tests in both mathematics and English. Students with lower academic or socio-economic backgrounds did not increase their chances of receiving a top grade. My results indicate that teachers predicted grades did not change the ordering of students final grade across students with varying baseline chances of receiving a top grade, but actually amplifying the inequality among students in the same class rooms to receive a top grade. 

Furthermore, I find that gender affected the chances of receiving a high grade after controlling for the students' previous academic performance. Female students were more likely to receive top grades than male students under the non-standardized grading regime. Among students with a baseline probability 50\% of receiving an A or A*, male students were 6 percentage points less likely to receive the highest grades than their female counterparts. This estimate remains robust to the inclusion of subject-level controls, suggesting that female students received higher grades than male students with comparable prior academic performance within the same classroom. In contrast, I find only weak differences in outcomes between white and non-white students: white students were awarded top grades at a rate 1 percentage point higher than non-white students, but this difference becomes statistically insignificant once school-level controls are included.

My structural model predicts that the income inequality across students across university worsens after the shift from standardized exam to non-standardized exam. The linear coefficient between the average parental income scores and the selectivity recorded of a university increases more than 5 times after replacement. The coefficient for the share of female students and the selectivity of the university decreases from -0.07 to -0.25, and the coefficient for the average math score decreases from -1.0 to -1.6. My estimate predicts how the nonstandardized grades lead to wealthier students to move up the university rank and replacing students with less economic affluence. 

\vspace{0.5em}

My results describe how the removal of standardized exams led to a unification of achievement levels inputs across schools. Low quality schools took advantage of the grading discretion by employing a looser grading policy than high quality schools. In the resulting grade distribution, the gap in probabilities for the median student to receive a top grade in A-levels was reduced by [] percentage points. This aligns with \cite{figlio2004}'s findings of stronger concentration of grade manipulation at lower quality schools. 

I find that academically advantaged students increase their chances of receiving a top grade more than students with lower academic standings. This result aligns with the findings in \cite{TYNER2020102037}, that grade inflation tends to be larger for students with higher academic standings, regardless of their school of attendance. Although this suggests teachers or schools do not ignore the previous academic standings of their students, and potentially preserving the order of students in their classrooms by academic standings. However, my results suggest demographic backgrounds of students do contribute to differential grade inflation, as female students received higher grades than male students with similar academic standings. This aligns with \cite{lavy2018}'s findings, and suggests that discretionary grading benefits girls more as girls perform better in the daily classroom than boys, while under performing in high-pressure settings, such as the exam date. I do not find strong evidence of grade inflation by socioeconomic standings measured by the income deprivation index or regional proficiency to enter higher education. The findings are at odds with \cite{hanna2012} and \cite{burgess2013} which finds a positive relationship between socio-economic backgrounds and the extent of grade inflation.

My findings highlights the patterns of grade inflation dynamics within an institution. The extent of grade inflation is diverse within schools across multiple subjects, with their extent influenced by the the quantitative content of the subject as well as the teachers discretion towards inflating grades. Inflation levels exhibit upward persistence, with high inflating schools continuing to inflate the similar extent as the previous year, but lower inflating schools upward revising their inflation levels in the consecutive year. The persistence in grading standards are consistent with the findings in \cite{denning2023} which they find persistent loosening of grading standards at the individual level. 

Although my results indicate a larger shift in grades for students in socio-economically disadvantaged backgrounds, the improvements in grades do not necessarily translate to an increase in an improvement in placements. This is due to socio-economically disadvantaged students often do not apply to top ranked universities albeit receiving a high grade in A-levels. On the otherhand, students attending middle quality schools with good academic and socio-economic background leverage the upward shift in grades to secure placements at better universities. The inelasticity between grades and application to selective universities is similar to the results documented by \cite{NBERw18586}.  The upward shift of students from middle quality schools pulls students attending higher quality schools but with lower academic competence to lower their university placement. Despite positive relation between previous academic performance and the grade shift, the average academic score, measured by GCSE math, reduces at higher ranked universities. This indicates how non-academically related attributes, such as gender and socio-economic background dominates the upwards shifts due to academic performance. 

This paper contributes to the discussion on test-optional policies for university admissions. The results indicate how a removal of standardized exams changes the allocation of grades across the student population. 